% Generated draft: requires mathematical and source review.
\documentclass[11pt]{article}
\usepackage[margin=25mm]{geometry}
\usepackage{fontspec}
\setmainfont{Latin Modern Roman}
\usepackage{amsmath,amssymb,mathtools,unicode-math}
\setmathfont{Latin Modern Math}
\usepackage{enumitem,needspace,xurl,hyperref}
\setlength{\parindent}{0pt}
\setlength{\parskip}{5pt}
\setlength{\emergencystretch}{3em}
\newcommand{\N}{\mathbb{N}}
\newcommand{\Z}{\mathbb{Z}}
\newcommand{\Q}{\mathbb{Q}}
\newcommand{\R}{\mathbb{R}}
\newcommand{\C}{\mathbb{C}}
\newcommand{\F}{\mathbb{F}}
\newcommand{\A}{\mathbb{A}}
\newcommand{\PP}{\mathbb{P}}
\newcommand{\E}{\mathbb{E}}
\newcommand{\eps}{\varepsilon}
\newcommand{\dd}{\,\mathrm{d}}
\begin{document}
\section*{UnsolvedMath ID: 11000100}
\textbf{Source problem code:} catalogue entry
\subsection{Definitions and mathematical statement}
For the surface \(S\) intended by the source, ask whether
\[
 \forall\Gamma\leq\operatorname{Mod}_S,
 \quad[\operatorname{Mod}_S:\Gamma]<\infty
 \ \Longrightarrow\ H^1(\Gamma)=0.
\]
The coefficient group is not printed. In the usual characteristic-zero, trivial-coefficient interpretation, this is \(H^1(\Gamma;\Q)=\operatorname{Hom}(\Gamma,\Q)=0\), equivalently vanishing first Betti number. For finitely generated groups it means finite abelianization, not necessarily zero integral first homology.

The source supplies genus-two counterexamples, so the question cannot literally assert this for every surface type. It reports F. Taherkhani's 1999 MSU thesis computations giving finite-index genus-two subgroups with nonzero first cohomology, whereas all examined genus-three subgroups had zero first cohomology; it describes higher-genus calculations as beyond the resources then available. Those finite computations are not a proof for all genus-three subgroups. The intended remaining closed-surface range is suggested to be genus at least \(3\), but is not explicit in the first sentence.

It also reports J. D. McCarthy [McC2]: if \(S\) is closed of genus \(g\geq3\), and a finite-index subgroup \(\Gamma\) contains the Torelli group, then \(H^1(\Gamma)=0\). The Torelli group is the kernel of the action on integral first homology. The stated ingredients are Johnson's Torelli results, the congruence subgroup problem for \(\operatorname{Sp}_{2g}(\Z)\), and its Kazhdan property (T), in the printed range \(g\geq3\). This restricted result is not extended to subgroups that do not contain Torelli. The source also says \(H^1(\operatorname{Mod}_S)=0\), subject to the same unstated surface/coefficient context.
\subsection{Short English statement}
Must every finite-index mapping-class subgroup in the intended higher-genus setting have vanishing first Betti number? Keep the genus-two counterexamples, limited genus-three computations and Torelli-containing theorem separate.
\subsection{Sources}
\begin{itemize}
\item \emph{Problems on Mapping Class Groups and Related Topics}, edited by Farb, Question I, supplied PDF page 85. \url{https://www.math.uchicago.edu/~farb/papers/mcgbook.pdf}. Recorded access 2026-07-29; extraction pdf-font-structure. Page disambiguates the repeated question label.
\item F. Taherkhani, 1999 MSU Ph.D. thesis [Tah], source for the reported calculations; J. D. McCarthy [McC2], the Torelli-containing subgroup result. Full titles/bibliography entries and the Johnson/congruence/property-(T) references are not supplied.
\item Local UnsolvedMath registry, \texttt{data/unsolvedmath/problems.json}, permanent ID \(11000100\), exact source code \texttt{catalogue entry}; supplied OPEN-TRIAGE summary reports the general finite-index question as open.
\end{itemize}
\end{document}
